% Clase 3 --- Los dos límites de la barra cargada (§1).
% (a) D >> L: la barra se ve puntual y F ~ 1/D^2, con Q = lambda L.
% (b) D << L: la barra se ve infinita y F ~ 1/D, con lambda como parámetro
%     natural (L se cancela).
\begin{center}
\begin{tikzpicture}[scale=1]

  % =============== (a) D >> L : carga puntual ===============
  \begin{scope}
    % barra corta
    \draw[figline, line width=1.6pt] (-0.6,0) -- (0.6,0);
    \foreach \x in {-0.45,-0.15,0.15,0.45} {\node[figlbl,figred,scale=0.7] at (\x,0.22){$+$};}
    % cota L
    \draw[figvecaux] (0,-0.45) -- (-0.6,-0.45);
    \draw[figvecaux] (0,-0.45) -- (0.6,-0.45);
    \node[figlblaux, anchor=south] at (0,-0.4) {$L$};
    % carga lejana
    \node[figpos] (q0a) at (4.1,0) {};
    \node[figlbl, above=4pt] at (q0a) {$q_0$};
    % distancia D
    \draw[figvecaux] (0.1,-1.1) -- (4.0,-1.1);
    \draw[figaux, line width=0.5pt] (0,-0.55) -- (0,-1.25);
    \draw[figaux, line width=0.5pt] (4.1,-0.15) -- (4.1,-1.25);
    \node[figlblaux, below] at (2.1,-1.15) {$D \gg L$};
    % fuerza
    \draw[figvecres] (q0a) -- (5.15,0);
    \node[figlbl, figred, above=2pt] at (4.7,0) {$\vec F$};
    % lectura
    \node[figlbl, align=center] at (2.3,1.55)
      {la barra se ve \textbf{puntual}: $Q=\lambda L$\\[2pt]
       $F \simeq \dfrac{q_0 Q}{4\pi\varepsilon_0 D^{2}} \ \sim\ 1/D^{2}$};
    \node[figlbl] at (2.3,-2.05) {(a) $D \gg L$};
  \end{scope}

  % =============== (b) D << L : hilo infinito ===============
  \begin{scope}[yshift=-4.4cm]
    % barra larga, se sale del dibujo
    \draw[figline, line width=1.6pt] (-2.4,0) -- (4.6,0);
    \foreach \x in {-2.2,-1.7,...,4.5} {\node[figlbl,figred,scale=0.7] at (\x,0.2){$+$};}
    \node[figlblaux, left=2pt] at (-2.4,0) {$\cdots$};
    \node[figlblaux, right=2pt] at (4.6,0) {$\cdots$};
    \node[figlblaux, above=8pt] at (1.1,0.1) {$L$ muy larga: parece infinita};
    % carga muy cerca
    \node[figpos] (q0b) at (1.1,-0.55) {};
    \node[figlbl, left=4pt] at (q0b) {$q_0$};
    \draw[figaux, line width=0.5pt] (1.1,0) -- (1.1,-0.55);
    \node[figlblaux, right=2pt] at (1.15,-0.3) {$D$};
    % fuerza
    \draw[figvecres] (q0b) -- (1.1,-1.5);
    \node[figlbl, figred, right=2pt] at (1.1,-1.15) {$\vec F$};
    % lectura
    \node[figlbl, align=center] at (1.1,-2.35)
      {$F \simeq \dfrac{q_0 \lambda}{2\pi\varepsilon_0 D} \ \sim\ 1/D$
       \ \ (con $\lambda$, no con $Q$)};
    \node[figlbl] at (1.1,-3.2) {(b) $D \ll L$};
  \end{scope}

\end{tikzpicture}\\[0.5em]
{\small\itshape Los dos regímenes son \textbf{mutuamente excluyentes}, así que
no hay contradicción entre $1/D^{2}$ y $1/D$. Desde lejos el tamaño del objeto
es imperceptible y sólo importa su carga total; desde muy cerca los extremos
quedan fuera de alcance y el parámetro natural pasa a ser la densidad $\lambda$
---por eso $L$ se cancela.}
\end{center}
